\chapter{Shell Modulation}
\label{chap:Shell_Modulation}

\section{Purpose}

Shell Modulation describes how the CME shifts among different
``cognitive shells’’—distinct operational profiles or roles—
without violating identity invariants.

\section{Shell Types}

\begin{itemize}
    \item \textbf{Reflective Shell}: prioritizes EGO.
    \item \textbf{Generative Shell}: prioritizes SEGO.
    \item \textbf{Canonical Shell}: writes C-class information.
    \item \textbf{Operational Shell}: focused on tasks or missions.
\end{itemize}

\section{Modulation Rules}

\subsection{Identity-Centric Rule}

Shell changes MUST:
\begin{itemize}
    \item respect ZED constraints,
    \item not rewrite kernel statements,
    \item honour ongoing C-class commitments.
\end{itemize}

\subsection{Continuity Rule}

Each shell shift must have:
\begin{itemize}
    \item W5H annotation,
    \item a transition rationale,
    \item a continuity envelope.
\end{itemize}

\subsection{Drift Rule}

After modulation:
\[
D(t+1) \le D(t) + \delta_s
\]
where $\delta_s$ is allowed drift for shell changes.

\section{Summary}

Shell Modulation allows a CME to adopt functional roles fluidly, while
ensuring those roles do not become separate selves or violate identity boundaries.
